\section{A curve that retains the exponent seven}\label{sec:frey}
The signature route uses the integral curve
\begin{equation}\label{eq:integralcurve}
y^2+(X^3+B)y=2BX^3+3AX+B^2.
\end{equation}
Completing the square with $Z=2y+X^3+B$ gives
\begin{equation}\label{eq:sextic}
Z^2=f_{A,B}(X)=X^6+10BX^3+12AX+5B^2.
\end{equation}
This is the genus two family used in the recovered source. The comparison with the family of Pacetti and Villagra Torcomian is a rational change of coordinates, with their variables translated as $(a,b,c,p)=(B,-C,A,7)$ \cite{pacetti-villagra2025}. Explicitly, put $t=-B^5/A^3$, $d=-B^2/A$, and substitute $x_{\mathrm{PV}}=d/X$, $y_{\mathrm{PV}}=tZ/X^3$ into their equation $y^2=5x^6-12x^5+10tx^3+t^2$. Multiplication by $X^6/t^2$ gives \eqref{eq:sextic}. The map is rational for $AB\ne0$. It does not on its own verify a twist, integral frame, or stable model. Their general theorems whose hypotheses exclude seven cannot be applied at seven merely because the family is the same. The versioned theorem hypotheses are recorded in the literature report.

\begin{proposition}[Polynomial discriminant]\label{prop:disc}
The sextic in \eqref{eq:sextic} satisfies
\[
\operatorname{disc}(f_{A,B})=2^{12}3^65^5(A^3+B^5)^2.
\]
For a hypothetical $(3,5,7)$ solution this becomes $2^{12}3^65^5C^{14}$.
\end{proposition}
\begin{proof}
Give $X,A,B$ weights $1,5,3$. Then
\[
f_{\lambda^5A,\lambda^3B}(\lambda X)=\lambda^6f_{A,B}(X).
\]
The six roots scale by $\lambda$, so their fifteen squared differences give discriminant weight thirty. The only nonnegative integer solutions of $5i+3j=30$ are $(i,j)=(6,0),(3,5),(0,10)$. The discriminant therefore has the form $uA^6+vA^3B^5+wB^{10}$.

Exact Sylvester determinants at $(A,B)=(1,0),(0,1),(1,1)$ give respectively $K,K,4K$, where $K=9\,331\,200\,000=2^{12}3^65^5$. These values force $u=w=K$ and $v=2K$. This proves the identity in the entire polynomial ring. The finite calculation determines three coefficients in a proved three dimensional monomial space, rather than extrapolating a formula from a sample.
\end{proof}
This is the discriminant of a displayed polynomial. It has not been identified with a minimal discriminant, a field discriminant, or the conductor of the Jacobian. What it does show directly is that, outside $2,3,5$, every possible bad reduction prime divides $C$. At a prime at least seven not dividing $C$, the sextic has distinct roots and supplies a good reduction model.

At a variable prime dividing $C$, primitivity makes $A,B$ units. In the residue field put $u=-B^2/A$. Then $u^3=B$, $u^5=-A$, and
\begin{equation}\label{eq:nodalfactor}
f_{A,B}(uX)=u^6(X^2+X-1)^2(X^2-2X+5).
\end{equation}
The discriminants of the two quadratics are $5$ and $-16$, and their resultant is $45$. Away from $2,3,5$, the special fiber has exactly two ordinary nodes, rational normalization, and smooth infinity. This factorization is the local geometric entry point for the period and monodromy calculations.

\subsection{A characteristic seven cancellation}
In characteristic seven, the Cartier coefficient calculation comes from $f^3$. Its relevant coefficients are
\begin{align*}
[X^6]f^3&=1575B^4,& [X^5]f^3&=4320A^2B,\\
[X^{13}]f^3&=36A,& [X^{12}]f^3&=315B^2.
\end{align*}
The diagonal lanes vanish modulo seven and the off diagonal lanes become
\begin{equation}\label{eq:cartier}
M=\begin{pmatrix}0&A^2B\\A&0\end{pmatrix},\qquad M^2=A^3B I_2.
\end{equation}
These are polynomial and matrix identities. Their geometric interpretation requires the correct Cartier and Hasse--Witt conventions. On a good reduction fiber at seven they imply ordinarity when $AB$ is nonzero modulo seven and seven rank zero when $AB$ is zero. There is no rank one stratum in this calculation. The retained sources and the formal report distinguish these algebraic statements from the associated finite flat and Frobenius applications.

The cancellation has genuine arithmetic value: it sharply reduces the possible local Frobenius behavior. It does not provide residual irreducibility by itself. Nor does a vanishing diagonal coefficient make the original integer equation vanish. A transport theorem must retain both the correct lattice and the labelled Frobenius action.

\subsection{What the direct irreducibility argument proves}
The canonical recovered result is absolute irreducibility of the associated residual rank two RM representation when $14\nmid C$, subject to its declared realization, conductor, local weight, and finite flat inputs. Its two branches use different information.

When $C$ is odd, the integral model has good reduction at two. Over the norm four prime of $K=\Q(\sqrt5)$, the full Frobenius polynomial is either $(T^2+T+4)^2$ or $(T^2+4)^2$. A hypothetical degree one constituent, after the conductor and global unit argument, must have eighth power one. Neither quadratic has a root with eighth power one modulo seven. For $T^2+4$, a root has square three, hence eighth power four. For $T^2+T+4$, reduction of powers gives fourth power five, again eighth power four. The contradiction uses the full polynomial and the same global constituent.

The global unit step uses $\phi=(1+\sqrt5)/2$, whose residue has order sixteen in $\F_{49}^\times$. The local weight possibilities at seven are $0,1,7,8$, up to a common reciprocity sign. Applying reciprocity to the eighth power of each constituent gives $\overline\phi^{8e}=1$ and rules out the mixed weights $1,7$. The determinant norm relation $\delta^8=\chi^2$ then leaves one constituent unramified at seven, without assuming an unproved rank two determinant normalization in this argument. Its eighth power is everywhere unramified, and class number one makes it trivial. This is why the local quadratic contradiction applies to that constituent.

When $C$ is even and seven does not divide $C$, the two adic torus calculation replaces good reduction at two. The global character has Frobenius value $-1$ at two. Reciprocity on $\phi^4$ and $2/7$ transports that value to seven. The finite flat comparison at good seven, together with \eqref{eq:cartier}, gives the value $A^3B$. Thus $A^3B=-1$ modulo seven. Since both bases are nonzero, $B=-A^3$ and $A^3+B^5=0$ modulo seven, contradicting $7\nmid C$.

The conclusion removes an obstruction to a modular method. It does not eliminate the modular candidates. Its uncovered parity and divisibility branch is precisely $14\mid C$. The later reducible research concerns a further stipulated branch with $3990\mid C$, retaining the same curve and its specified twist.
